\section{P2: give sparse-list expansion one monomial-sum view}

\subsection{The repeated reasoning}

\src{SyntaxCoefficients}{the interpretation lemmas} keeps sparse expansions
as executable lists. A term consists of an exponent vector and an integer
coefficient. Multiplication distributes through \code{flatMap} and
\code{map}, and duplicate monomials remain in the list.

The recursive \code{interpret} function denotes the corresponding ordinary
polynomial sum. The original proof layer independently inducts over lists
for append, coefficient negation, multiplication by one term, full
multiplication, and coefficient extraction. These are all algebraic facts
about list sums.

Keep the representation and every executable definition. Add just one
private view:
\[
 \operatorname{interpret}(L)=
 \sum_{t\in L}\operatorname{monomial}(\operatorname{exponent}(t_1),t_2),
 \tag{5}
\]
where the sum is an occurrence-sensitive list sum, not a finite set sum.

\subsection{The new proof structure}

Prove (5) once by list induction. Then use the following interfaces:

\begin{itemize}
\item Append: \code{List.map_append} and \code{List.sum_append}.
\item Negation: the additive-map law for monomial coefficients and
\code{List.sum_neg}.
\item Multiplication: flatten a list of lists, then use
\code{List.map_flatten}, \code{List.sum_flatten}, and distribution over sums.
\item Coefficient extraction: apply the bundled
\code{MvPolynomial.coeffAddMonoidHom} to the list sum.
\end{itemize}

The complete multiplication proof becomes
\begin{Lean}
  simp only [interpret_eq_sum, multiply, List.flatMap_def,
    List.map_flatten, List.sum_flatten, List.map_map,
    Function.comp_def, exponent_add,
    ← MvPolynomial.monomial_mul,
    List.sum_map_mul_left, List.sum_map_mul_right]
\end{Lean}

The private left-multiplication induction is no longer needed. The structural
induction proving \code{interpret_expand} remains: it is the correct connection
between expression constructors and the established list algebra.

A small companion derives \code{exponent_inj} from the injectivity of
\code{Finsupp.equivFunOnFinite.symm}. It replaces a local extensionality proof,
and is counted within this module change rather than as a separate proposal.

\subsection{Why the executable behavior stays exact}

Replacing the list by a finite set would discard duplicate occurrences.
This refactor does not do that. Empty lists, repeated exponent vectors,
negative coefficients, and coefficients that later cancel are still processed
by the original list operations.

Even the existing recursive \code{interpret} definition is retained, although
it is used only for mathematical interpretation. Equation (5) is a theorem
about it, so its definitional behavior is preserved.

\subsection{Measurements and checking scope}

The production candidate is 110 physical lines instead of 120, and 87
nonblank lines instead of 97. These counts include the new private theorem
and the explicit list-algebra import. Five list inductions become one; the
expression induction remains unchanged.

Both complete modules compile with Lean 4.32.0, using a freshly rebuilt
current \code{Representation} module. Both public
\code{coefficient_expand} endpoints report
\[
 \{\text{\texttt{propext}},\ \text{\texttt{Classical.choice}},\ \text{\texttt{Quot.sound}}\}.
\]
There is no downstream-build or speedup claim. The compiler times in the logs
are execution receipts from a shared environment, not benchmarks.

The exact patch is \code{patches/syntax-coefficients.patch}. Complete before
and after sources are also included.

\section{P3: transport quantification by restriction surjectivity}

\subsection{The actual environment-extension problem}

\src{ObjectPolynomial}{\code{sat_noZeroForm_iff}} relates validity of a sealed
no-zero formula over environments $\N\to V$ to a predicate over finite tuples
$\operatorname{Fin}(n)\to V$. Its original proof constructs an environment by
hand: use the tuple below $n$, and a default object elsewhere. It then proves
restriction recovers the tuple and transports guards and equation
satisfaction separately.

The content is simply surjectivity of restriction. For an injection
$j:I\hookrightarrow J$ and a default $d\in V$, the map
\[
 \rho:(J\to V)\longrightarrow(I\to V),\qquad
 \rho(e)=e\circ j
\]
is surjective. Thus every predicate $P$ on assignments satisfies
\[
       (\forall e,\ P(e\circ j))
       \quad\Longleftrightarrow\quad
       (\forall a,\ P(a)).
 \tag{6}
\]
Include the naturalness guard and unsolvability conclusion in $P$.

\subsection{Use the existing extension principle}

The checked generic mechanism uses Mathlib's
\code{Function.Injective.surjective_comp_right'} and
\code{Function.Surjective.forall}:

\begin{Lean}
theorem forall_assignments
    (coordinates : ι → I)
    (hinj : Function.Injective coordinates)
    (default : V) (P : (ι → V) → Prop) :
    (∀ env : I → V, P (env ∘ coordinates)) ↔
      ∀ assignment, P assignment :=
  (hinj.surjective_comp_right' (fun _ => default)).forall.symm
\end{Lean}

For the actual theorem, the injection is \code{Fin.val} and the default is
\code{vempty}~\code{H}. No new production helper is necessary: invoke the library
fact after the existing satisfaction equivalences.

\begin{Lean}
  refine (seal_valid (mem := mem) _).trans ?_
  change (∀ ee : Nat → V,
    Sat mem ee (inputsNaturalForm (Fin.val : Fin n → Nat)) →
      ¬ Sat mem ee (equationForm e Fin.val)) ↔ NoZeroRel H e
  simp only [inputsNaturalForm_spec H, equationForm_spec H]
  exact
    (Fin.val_injective.surjective_comp_right'
      (fun _ => vempty H)).forall.symm
\end{Lean}

The displayed source is wrapped for the page; the supplied patch contains
the exact eight-line declaration.

\subsection{Boundary cases and distinctness}

Injectivity matters. If two tuple coordinates map to the same environment
coordinate, arbitrary distinct tuple values cannot be recovered. The
universal assignment conclusion would then be stronger.

The explicit default handles the empty coordinate type without silently
assuming a nonempty target type in the abstract theorem. In the model
application, the empty set already provides this value. No standardness,
arithmetic, well-foundedness, or induction property of $V$ is used.

Earlier sentence-transport proposals consume the no-zero equivalence.
P3 instead simplifies the environment extension inside that equivalence,
so the benefit remains after those earlier wrappers have been applied.

\subsection{Evidence and outstanding integration}

The generic assignment theorem, its \code{Fin.val} specialization, and the
guarded-unsolvability specialization compile. Their only audited axiom
dependencies are the usual \code{propext}, \code{Classical.choice}, and
\code{Quot.sound}.

The production declaration changes from 20 to 8 lines, a proposed
twelve-line reduction with no added helper. However, the complete
\code{ObjectPolynomial} module was not compiled: its deep object-model
dependency closure was unavailable in the local cache. The production patch
is marked \emph{integration pending}. Success of the abstract mechanism is
not presented as a complete object-module check.

\section{P4: use the order obstruction on a residual exponent}

\subsection{A smaller proof of the existing theorem}

\src{NativeBinaryScale}{\code{width_dvd_exponent}} proves
\[
      2^d-1\mid 2^u-1\quad\Longrightarrow\quad d\mid u
\]
with its retained public assumption $0<d$. Put $r=u\bmod d$. The existing
modular-power identity reduces the divisibility hypothesis to
\[
      2^r-1=0.
\]
The old proof turns this into $2^r=1$, rewrites $1$ as $2^0$, and invokes
injectivity of exponentiation.

The replacement uses the order obstruction:
\[
      2^r-1=0\Longrightarrow 2^r\le1.
\]
If $r\ne0$, \code{Nat.pow_le_one_iff} would force $2\le1$.
So $r=0$, and hence $d\mid u$.

\begin{Lean}
  have hm := Nat.mod_eq_zero_of_dvd hdiv
  rw [Nat.pow_sub_one_mod_pow_sub_one] at hm
  apply Nat.dvd_of_mod_eq_zero
  by_contra hn
  have hsmall :=
    (Nat.pow_le_one_iff hn).mp (Nat.sub_eq_zero_iff_le.mp hm)
  omega
\end{Lean}

The public signature is unchanged, including \code{hd}. The supplied source
formats the body in six lines; the page breaks one long line.
The declaration shrinks from twelve to eight lines. Both exact declarations
compile with identical axiom sets. Their unused-\code{hd} warning is also
identical.

\subsection{Optional arithmetic API: recover a block power directly}

The independently checked \code{PrimeRadix.lean} proves
\[
\begin{gathered}
 p\text{ prime},\quad 1<q,\quad q\mid p^R,\quad p^d-1\mid q-1\\
 \Longrightarrow\quad
 \exists N>0,\quad dN\le R,\quad q=(p^d)^N .
\end{gathered}
\tag{7}
\]
Classify a divisor of a prime power as $q=p^u$ with $u\le R$.
The generalized exponent-divisibility fact yields $u=dN$, and $q>1$
forces $N>0$.

An affine adapter accepts $q=1+J(p^d-1)$, matching \code{sourceQ}.
A source client can then supply its divisibility certificate and recover
the cell-count exponent and budget together, rather than exposing an
intermediate binary exponent.

The generic core needs no separate $d>0$ hypothesis: for $d=0$ the other
premises force $q=1$, contradicting $q>1$. The public source wrapper can
still keep its original parameter.

This optional API is not counted as additional removed lines. Its shared
proof cost may exceed its benefit for a single current client. Adopt it
if the arithmetic classification becomes a reused boundary.

\subsection{Necessary guards}

The $q>1$ premise is required for a positive exponent. Primality is also
substantive. For
\[
       p=6,\quad d=2,\quad q=4096,\quad R=12,
\]
one has $4096\mid6^{12}$ and $35\mid4095$, but $4096$ is not a positive power
of $36$. Thus (7) cannot be generalized to arbitrary bases larger than one.
Exponent divisibility alone needs only $p>1$; primality is needed to
classify divisors of $p^R$.

The exact declaration checks and the optional general prototype are
included. The modified complete \code{NativeBinaryScale} module and its
downstream clients were not rebuilt.
